\section{Related Works}
\label{sec:related_work}
\subsection{Graph Representation Learning}
Graph representation learning aims to capture structural and semantic information from graph-structured data. GNNs have become a dominant tool in this area by aggregating information through message passing~\cite{cai2018comp,liu2021pick,DBLP:conf/nips/BodnarFOWLMB21,bodnar2021weisfeiler}. Despite their success, GNNs struggle with dynamic or multi-task settings due to rigid architectures and fixed message-passing schemes that lack task-aware modulation and flexible inference control.%Despite their success across diverse tasks, GNNs often struggle to adapt to changing objectives due to their rigid architectures and fixed message-passing schemes. These models typically lack mechanisms for task-aware modulation or flexible control during inference, which limits their application in dynamic or multi-task settings.

To improve model flexibility and task adaptivity, recent studies have proposed graph prompt learning, inspired by prompt-based tuning in natural language processing~\cite{wei2023chainofthought,zhou2023docprompting}. This paradigm typically involves pre-training GNNs on large-scale graph~\cite{hu2020gpt,graphcl,qiu2020gcc,zhao2024all,yu2023generalized} and guiding downstream task execution through prompt-like structures that inject task-relevant information.
For example, VNT~\cite{VNT} introduces virtual nodes as soft prompts to encode task-level context and guide message propagation. GraphPrompt~\cite{graphprompt} proposes a task-specific readout mechanism that learns to aggregate node representations differently depending on the downstream objective. GraphPro~\cite{yang2024graphpro} further designs spatial- and temporal-aware gating mechanisms, adapting the prompt to dynamic graph settings like recommender systems. PRODIGY~\cite{mishchenko2023prodigy} constructs a task graph as an explicit prompt structure and learns cross-graph in-context learning capabilities by interacting it with a data graph via a transformer-based module. ProNoG~\cite{Yu2025ProNoG} proposes a novel pre-training and prompt learning framework tailored for non-homophilic graphs, introducing a conditional network that captures node-specific relational patterns to enhance downstream task performance. 
Graph prompt learning typically presumes structural alignment with pretraining data, hindering generalization to graphs with divergent topologies.
%These methods enable few-shot generalization and reduce the need for task-specific re-training by conditioning the model on external task instructions or auxiliary graph structures. 
%However, despite improving task adaptation, graph prompt learning typically presumes alignment between input and pretrained structural distributions, leading to poor generalization on structurally divergent graphs due to the absence of external structural knowledge.
%To improve model flexibility and task adaptivity, recent studies have proposed graph prompt learning, inspired by prompt-based tuning in natural language processing~\cite{wei2023chainofthought,zhou2023docprompting}. This paradigm involves pre-training GNNs on large-scale graph data—either labeled or unlabeled~\cite{hu2020gpt,graphcl,mishchenko2023prodigy,zhao2024all,xia2024opengraph}—and guiding downstream task execution via prompts, such as virtual nodes~\cite{VNT}, task-aware readout heads~\cite{graphprompt}, or task graphs~\cite{mishchenko2023prodigy}. These techniques enable in-context learning (ICL) and few-shot generalization across multiple tasks, reducing the need for full model fine-tuning. 
%However, while graph prompt learning enhances task adaptation, it generally assumes that the structural distribution of input graphs aligns with that seen during pretraining. When applied to unseen graphs with significantly different structural patterns, these methods still struggle due to the lack of access to external structural knowledge.

To overcome this limitation, recent work has drawn inspiration from RAG in NLP~\cite{lewis2021retrievalaugmented,zhao2023retrieving}, where external information is retrieved and injected into a pre-trained model to support better generalization. In the context of graph learning, a representative effort is RAGRAPH~\cite{JiangQXZZ00X0W24}, which enhances GNNs through a retrieval-based framework that leverages a toy-graph library to incorporate structurally or semantically similar graphs via message-passing prompting, enabling adaptation without retraining. %which proposes a retrieval-based framework to enhance GNNs by leveraging a curated library of toy graphs. During inference, RAGRAPH retrieves structurally or semantically similar graphs from this library and incorporates them into the input via a message-passing prompting mechanism, enabling the model to adapt its behavior based on retrieved structural patterns without retraining. 
Despite their effectiveness, existing graph RAG methods primarily rely on low-dimensional structures, overlooking higher-dimensional topologies crucial for capturing key substructures and complex relations.

% Additional related works on topological deep learning are discussed in Appendix~\ref{app:related}.

\subsection{Topological Deep Learning} 
Topological deep learning (TDL) builds upon early advances in Topological Signal Processing (TSP) \cite{barbarossa2020topological, schaub2021signal, roddenberry2022cellsp, sardellitti2022cell}, which highlighted the importance of modeling higher-order relations beyond pairwise connections. This line of work has motivated the generalization of classical graph-theoretic tools to richer topological domains, such as simplicial complexes (SCs) and cell complexes (CWs). In particular, extensions of the Weisfeiler–Lehman test to SCs and CWs \cite{bodnar2021weisfeiler, DBLP:conf/nips/BodnarFOWLMB21} have laid the theoretical foundations for higher-order message passing. Subsequently, a variety of neural architectures have been proposed, including convolutional designs for SCs and CWs \cite{ebli2020simplicial, yang2022effsimpl, hajij2020cell, yang2023convolutional, roddenberry2021principled, hajij2022} and attentional formulations \cite{goh2022simplicial, giusti2023cell}. To unify these efforts, the combinatorial complex (CC) framework \cite{hajij2023topological} was introduced, providing a general message-passing paradigm that encompasses SCs, CWs, and hypergraphs. Parallel to these developments, Sheaf Neural Networks (SNNs) \cite{HGh19, hansen2020sheaf, sheaf2022, battiloro2022tangent, battiloro2024tangent, barbero2022sheaf} have demonstrated the effectiveness of sheaf-based modeling in handling heterophily and capturing localized topological constraints.

While these works substantially extend graph representation learning to higher-dimensional domains, they largely remain confined to direct message passing within predefined complexes. Such approaches often exhibit limited generalization to unseen graphs with substantially different topologies, a limitation that poses significant challenges for real-world applications where structural variability is the norm. %In contrast, our work goes beyond intra-complex learning by \textit{retrieving and integrating multi-dimensional cellular topologies} from external corpora, enabling topology-aware retrieval and multi-level reasoning that enhance adaptability across diverse graph scenarios.